% Created 2026-06-08 Mon 16:11
% Intended LaTeX compiler: pdflatex
\documentclass[11pt]{article}
\usepackage[utf8]{inputenc}
\usepackage[T1]{fontenc}
\usepackage{graphicx}
\usepackage{longtable}
\usepackage{wrapfig}
\usepackage{rotating}
\usepackage[normalem]{ulem}
\usepackage{amsmath}
\usepackage{amssymb}
\usepackage{capt-of}
\usepackage{hyperref}
\author{fcalado}
\date{\today}
\title{}
\hypersetup{
 pdfauthor={fcalado},
 pdftitle={},
 pdfkeywords={},
 pdfsubject={},
 pdfcreator={Emacs 29.3 (Org mode 9.6.15)}, 
 pdflang={English}}
\begin{document}

\tableofcontents

\section{reader reports}
\label{sec:orgb69919d}
\subsection{feedback}
\label{sec:org2cb203e}
\begin{itemize}
\item overall issues, how the project is currently constituted
\begin{itemize}
\item how does \textbf{argument accumulate} in each chapter?
\begin{itemize}
\item connection between object and theory ought to feel necessary
\item stronger grounding in book's overall argument
\item "what the bringing together of queer studies/theory with digital
humanities tools and technologies does to transform or critique
the field"
\begin{itemize}
\item "what takeaway does this give us that is broadly applicable
within that field? What does it reveal about either the
technology discussed or queer theory broadly?
\end{itemize}
\end{itemize}
\item where is the \textbf{urgency}?
\begin{itemize}
\item feels like series of term papers
\item needs strong contribution to the discipline
\end{itemize}
\item what is the \textbf{methodological} intervention?
\begin{itemize}
\item case needs to be made for the innovativeness of methodology, as
it appears to be mostly close reading for now
\begin{itemize}
\item how does cllose reading work wrt digital tools?
\begin{itemize}
\item "its reliance on close reading seems to work against its
goal of working as an introduction to digital tools"
\end{itemize}
\end{itemize}
\item consider \emph{Data Feminism} as a model
\end{itemize}
\end{itemize}
\item length can be expanded
\begin{itemize}
\item currently between 7000-8000 words
\item more space devoted to the political or theoretical intervention
of each chapter
\item the ambition and length of the chapters expanded
\end{itemize}
\item citing prior Dh work
\begin{itemize}
\item Jason Boyd, Bonni Ruberg, Rahul K Gairola, Ashley Caranto,
attendees of first Queering the Digital Humanities course, DHSI in
Victoria, Canada, 2018. All four have published substantial work
on the topic since that first session.
\item Also look earlier – DH emerged from cultural studies, Stuart
Hall’s “Encoding/ Decoding” to give a rich and balanced,
historical literature review.
\item adding previously-published work that has appeared since the
author completed the PhD thesis
\end{itemize}
\item title change
\begin{itemize}
\item Queer Revisions of Language Technology: Contemporary Counter-Codes
of Language Data
\end{itemize}
\item writing style
\begin{itemize}
\item ensure that the writing is as lucid and accessible as possible so
that a wide swathe of reads – not just data scholars – can access
the important literary, DH, and LGBTQIA+ readings.
\end{itemize}
\end{itemize}
\subsection{response}
\label{sec:org6a53707}
\subsubsection{How does the argument accumulate over the chapters?}
\label{sec:org56f90ed}
Language technologies reduce language complexity to make it
computable. Words are removed from their original context, cleaned and
standardized, so that they can be stored and analyzed at scale.

EITHER

These increasingly automated processes encode dominant logics and
extractive incentives onto language data, creating known issues
with data sovereignty, privacy, and most recently, bias, as AI
tools.

OR

As a result of this process, dominant forms of language meanings
are reproduced into digital realms, creating known issues like bias
and discrimination, as with current AI tools.

This book breaks open and re-works technologies for analyzing language
about sex, gender and sexuality. It takes a close look into
technologies like web scraping, text analysis, machine learning, among
others. Then, drawing from concepts in Queer Studies and related
fields, it re-works these technologies into new analytical
methodologies that open up, rather than reduce, novel readings of
language data.

Each of the five chapters considers a different technology for working
with language. For that technology, the chapter identifies a technical
constraint that reduces or flattens language meaning. For example,
chapter 2, "Text Analysis," explores the technology of quantitative
text analysis, and within that technology, the constraint of
iteration: a logic of repetition that allows computers to execute the
same action across large collections of text, one word at a time. This
constraint, which allows quantitative text analysis to transform texts
at large scales, amplifies frequent patterns in language while
supressing minor differences and details. The chapter then takes
Judith Butler's concept of Gender Performativity, which describes
gender as a social and behavioral phenomenon that becomes intelligible
through habit or repetition, to re-work computational iteration into a
new text analysis practice for studying gender. Unlike typical text
analysis projects, which visualizes dominant semantic patterns and
relationships, this new methodology uses quantification to surface the
idiosyncracies and nuances in word meaning.

The table below sketches the relationship between technology,
technical constraint, and conceptual tool across the book's five
chapters.

\begin{center}
\begin{tabular}{llll}
Chapter & Technology & Technical Constraint & Conceptual Tool\\[0pt]
\hline
1. Text Capture & Web Scraping & Protocol & Borderlands Theory\footnotemark\\[0pt]
2. Text Analysis & Quantitative Text Analysis & Iteration & Gender Performativity\footnotemark\\[0pt]
3. Text Preservation & Markup Language & Containment & Dis-Identification\footnotemark\\[0pt]
4. Text Display & Display & Abstraction & Enfleshment\footnotemark\\[0pt]
5. Text Generation & Machine Learning & Aggregation & Body Dysphoria\footnotemark\\[0pt]
\end{tabular}
\end{center}\footnotetext[1]{\label{org35d1606}Anzaldúa, Gloria. \emph{Borderlands/La Frontera: The New Mestiza}.
Aunt Lute Books. 1987.}\footnotetext[2]{\label{orgfce2e33}Butler, Judith. \emph{Bodies That Matter: On the Discursive Limits
of “Sex”}. Routledge, 2011.}\footnotetext[3]{\label{orgd4c360d}Muñoz, José Esteban. \emph{Disidentifications: Queers of Color and
the Performance of Politics}. José Esteban Muñoz. University of
Minnesota Press, 1999.}\footnotetext[4]{\label{org1c1542c}Musser, Amber Jamilla. \emph{Sensual Excess: Queer Femininity and
Brown Jouissance}. NYU Press, 2018.}\footnotetext[5]{\label{orgd984241}Prosser, Jay. \emph{Second Skins: The Body Narratives of
Transsexuality}. Columbia University Press. 1998.}


Each chapter tests its methodology on a different literary object.
While the earlier chapters explore traditional works of queer
literature (Gloria Anzaldúa's \emph{Borderlands/La Frontera: The New
Mestiza}, 1987; Virginia Woolf’s \emph{Orlando: A Biography}, 1928; Oscar
Wilde’s \emph{The Picture of Dorian Gray}, 1890) the latter chapters take
on more eclectic source material, some of which features heterosexual
characters, from science fiction, electronic literature and reality TV
(Octavia Butler's \emph{Dawn}, 1987; Entropy8Zuper!'s \emph{skinonskinonskin},
1999; \emph{Love is Blind}, 2020+). The expansion of source material points
new potential orientations around queerness based on embodiment and
mutual feeling, following Butler's famous call for queerness which is
"never fully owned, but always and only redeployed, twisted, queered
from a prior usage and in the direction of urgent and expanding
political purposes."\footnote{Butler, \emph{Bodies That Matter}, 173.} (173).

Part of re-defining "queer" for the future involves developing a
robust conceptual toolset for doing resistance work. As the earlier
chapters show, Queer Studies' initial (and necessary) emphasis on
identity politics offers resistance methods from \emph{within} systems of
power. The first three chapters thus explore ways for re-working
protocols and structures in order to resist totalizing logics.
However, the latter chapters, which draw from more recent Trans- and
Black Feminist-inflected work in Queer Studies, foreground
methodologies of resistance that work \emph{alongside} and \emph{beyond} these
systems. These methodologies center on material and embodied
experiences in ways that are not contained within existing systems.
The book argues that both approaches, working \emph{inside} and working
\emph{beyond}, are necessary for resistance, which must shift between
opposition and collaboration, from co-opting to co-creating within and
without oppressive forces.

\subsubsection{What is the urgency?}
\label{sec:org7bf1349}
Language technologies are increasingly centralized, black-boxed, and
automated. This book offers a blueprint for how language technologies
work from the inside. It examines a "stack" of technologies that
progress from preliminary to advanced. The book begins with three
chapters on the preliminary processes of gathering, cleaning, and
structuring text ("Text Capture," "Text Analysis," "Text
Preservation"). Then, the fourth and fifth chapters proceed to the
more advanced tasks of screen display and machine learning ("Text
Display," "Text Generation").

This scaffolded approach builds a knowledge base over a variety of
tools through the course of the book. Through exposure to the logics
of computational procedures, readers learn technical strategies to
re-work the all but inevitable language technologies that proliferate
around them.

\subsubsection{What is the methodological intervention?}
\label{sec:orgbca4e0f}
The main intervention is the close reading of digital tools, a
methodology that it adapts from Matthew Kirschenbaum.

"its reliance on close reading seems to work against its goal of
working as an introduction to digital tools"


The chapter then applies a concept from Queer Studies, which it uses
as a conceptual \emph{tool} to re-work this constraint to surface what has
been supressed.
\end{document}
